\textbf{DPLL} is an algorithm to establish the satisfiability or unsatisfiability of a CNF propositional formula. \vspace{7pt}

The basic idea of this algorithm is divided into three main steps:
\begin{enumerate}
    \item Apply \textbf{unit resolution} as long as possible. At some point we will not able to derive any other clause, so we integrate search processes, something similar we did with Constraint Programming.
    \item \textbf{Choose} a boolean decision variable $p$, introducing inside the search tree the cases $p$ and $\neg p$.
    \item After we introduced the positive and negative cases, we \textbf{go recursively} until we reach the satisfiability or unsatisfiability condition.
\end{enumerate}
\begin{definition}[title = DPLL pseudocode]
    \begin{algorithmic}
        \Procedure{DPLL}{$X$}
            \State $X = \text{unit-resol(X)}$
            \If{$X = \emptyset$}
                \State $\text{return}(\textit{sat})$
            \EndIf
            \If{$\bot \notin X$}
                \State $\text{choose variable } p \in X$
                \State $\text{DPLL}(X \cup \{p\})$
                \State $\text{DPLL}(X \cup \{\neg p\})$
            \EndIf
        \EndProcedure
    \end{algorithmic}

    Before moving on, let's take a look at the previous snippet of code. The main steps of the DPLL algorithm can be described as follows:
    \begin{enumerate}
        \item First of all, we apply \textbf{unit resolution} as long as possible.
        \item We check if we have reached the \textbf{contradiction}, if so we return the satisfiability condition.
        \item If we didn't come up with the bottom or the empty set, we add a new boolean variable inside the description of the problem.
        \item Finally, we augment the original set of propositional formulas, adding the positive and negative cases of boolean variable chosen.
    \end{enumerate}
\end{definition}
DPLL and Constraint Programming are very similar paradigms:
\begin{enumerate}
    \item Like Constraint Programming, we check during the resolution if we could reach a contradiction; if so, we backtrack choosing another variable or assigning another boolean value.
    \item Both of them are complete methods. Giving to them enough time, they always return a final answer that might be either positive or negative.
    \item Efficiency strongly depends on the choice of the variable.
\end{enumerate}
\begin{example}
    e.g. DPLL resolution process. \vspace{7pt}

    \begin{enumerate}
        \item $\neg p \lor \neg s$
        \item $\neg p \lor \neg r$
        \item $\neg q \lor \neg t$
        \item $p \lor r$
        \item $p \lor s$
        \item $r \lor t$
        \item $\neg s \lor t$
        \item $q \lor s$
        \item $q \lor \neg r$
    \end{enumerate} \vspace{7pt}

    This set of propositional formulas doesn't allow any unit resolution, since there is no boolean variable alone. Therefore, we need to choose a variable 
    and add it inside this set, defining its positive or negative value. \vspace{7pt}

    Assuming we choose the boolean variable $p$ and we add it inside the resolution in its positive form $p = T$, the mechanism goes as follows:
    \begin{enumerate}
        \item $\neg s \ (1)$ 
        \item $\neg r \ (2)$
        \item $q \ (\neg s, 8)$ 
        \item $t \ (\neg r, 6)$
        \item $\neg t \ (q, 3)$
        \item $\bot \ (t, \neg t)$
    \end{enumerate} \vspace{7pt}

    Note: after we add in the set $p = T$ then $(1)$ and $(2)$ propositional formulas of the original list must be exactly true, therefore both $\neg s$ and 
    $\neg r$ must be equal to false. \vspace{7pt}

    As we can observe, we have reached a contradiction, so the problem is unsatisfiable \textit{(UNSAT)}. Now we try with its counterpart, adding to 
    the original list of propositional formulas $\neg p$.
    \begin{enumerate}
        \item $r(4)$
        \item $s(5)$
        \item $q(r, 9)$
        \item $t(s, 7)$
        \item $\neg t(q, 3)$
        \item $\bot (t, \neg t)$
    \end{enumerate} \vspace{7pt}

    Again, we have reached a contradiction, therefore the whole problem is UNSAT.
\end{example}